% -------------------------------------------------------------------
\section{Every difference from Kiem, accounted for}
\label{sec:kiemledger}

Table~\ref{tab:kiemledger} closes the chapter's account. It lists every number
in which our design and Kiem's differ, says whether the comparison is valid as
published, names the cause of the difference and points to the measurement
that shows it. Where the cause is one design decision, the evidence is the step
of the ladder that changes only that decision. Where it is not, the table says
so.

\begin{table}[p]
\centering
\caption[Every difference from Kiem's design: cause and evidence]{Every
difference between DRHE-1 and Kiem's published design, with its cause and the
evidence for it. ``Valid'' says whether the two published numbers measure the
same thing. Replica figures are post-route on the \texttt{xcku5p}; the replica
is the fair yardstick for what our design decisions cost.}
\label{tab:kiemledger}
\scriptsize
\renewcommand{\arraystretch}{1.25}
\setlength{\tabcolsep}{3pt}
\begin{tabular}{@{}>{\raggedright}p{0.09\textwidth}>{\raggedleft}p{0.085\textwidth}>{\raggedleft}p{0.08\textwidth}>{\raggedleft}p{0.08\textwidth}>{\raggedright}p{0.075\textwidth}>{\raggedright}p{0.30\textwidth}>{\raggedright\arraybackslash}p{0.16\textwidth}@{}}
\toprule
Metric & DRHE-1 & Replica & Kiem & Valid & Cause of the difference & Evidence\\
\midrule
LUT & \HW{drheone}{LUT} & \HW{ablk}{LUT} & \num{45892} & yes &
DRHE-1 against the replica: floating-point operators (\HW{drheone}{fopLUTpct}\,\% of DRHE-1's LUTs), then the packer. The replica against Kiem: about \SI{5}{\percent} explained (Table~\ref{tab:kiemgap}). &
ladder steps 1, 2, 4; Table~\ref{tab:kiemgap}\\
FF & \HW{drheone}{FF} & \HW{ablk}{FF} & \num{9243} & yes &
Registers inside float operators (\HW{drheone}{fopFFpct}\,\%) and along a 58-stage pipeline. Kiem's extra over the replica: about \SI{45}{\percent} explained (speed grade, packer, depth), the rest not. &
ladder steps 4, 5; Table~\ref{tab:kiemgap}\\
DSP & 132 & \HW{ablk}{DSP} & 44 & yes &
Every DSP belongs to a float operator; each float multiply needs 3, each sine/cosine 13. Sharing at $\mathrm{II}=2$ halves the count. &
Table~\ref{tab:dspcount}; step~3 (\HW{ablptwor}{DSP} at $\mathrm{II}=1$)\\
BRAM (36\,Kb) & \HW{drheone}{BRAMT} & \HW{ablk}{BRAMT} & 10 & yes, in tiles &
32-bit state words need a full tile per array; 16-bit words need half. &
Table~\ref{tab:memtable}; replica matches exactly\\
II & 2 & 1 & 1 & yes &
One extra write per cycle to clear the state; three accesses on two ports. &
HLS~200-885; step~3 reaches II~1\\
Depth (cycles) & 58 & \HW{ablk}{hDepth} & 23 & yes, both HLS &
Closed loop: polar conversion waits for the prediction (26 cycles); the rest is float operator latency. & Figure~\ref{fig:sched}; open-loop float schedule (depth 32); steps 4, 5\\
$F_{\max}$ (MHz) & \HW{drheone}{fmax} & \HW{ablk}{fmax} & ${\geq}100$ & partly &
Both meet \SI{100}{\mega\hertz}; Kiem reports only that he met it. His part is speed grade $-1$. &
routed timing; $-1$ runs\\
Throughput (Gbit/s) & 6.40 th., 5.14 meas. & 12.80 th., \HW{ablk}{thrmeas} meas. & 12.80 th. & theoretical only &
II~2 against II~1. Kiem's 11.92 is binary units and theoretical. &
Eq.~5.1 of \cite{kiem2025}; co-simulation\\
Latency & 580\,ns pipeline & --- & 230\,ns pipeline; 0.08\,ms system & pipeline only &
Pipeline: the depth above. His 0.08\,ms spans FFT, DMA and interrupts. &
Section~\ref{sec:latency}\\
Power & \HW{drheone}{pdyn}\,W dyn. & \HW{ablk}{pdyn}\,W dyn. & not reported & --- &
Scales with switching resources; float design $\approx 2\times$ replica. &
Table~\ref{tab:power}\\
Compression ratio & 3.310 (50 frames; 3.321 on the first 10) & 3.319 (first 10 frames) & 3.17 & no &
Different sensor and data scaling; the float and fixed designs differ by \SI{0.07}{\percent} on the same data. &
Section~\ref{sec:crinterp}\\
Device fraction & 50\,\% of ZU3CG LUTs & 24\,\% & 65\,\% & yes &
Follows from the LUT row. &
Table~\ref{tab:resdevice}\\
\bottomrule
\end{tabular}
\end{table}

\paragraph{The story in one paragraph.} The published comparison said our
design was two and a half times larger than Kiem's, slower, and too big for his
device. None of that survives. The size and clock figures came from HLS
estimates set against his post-route figures; implemented, DRHE-1 uses fewer
LUTs than his IP, meets the same clock with a quarter of the period to spare, and
fits his device at half its logic. Against a replica of his design built in the
same tool, DRHE-1 is larger, and by exactly the amounts its four distinctive
decisions cost: floating-point arithmetic (most of the logic and registers,
every extra DSP, twice the block RAM), a less economical packer (about
\num{4200} LUTs), a closed loop (26 cycles of depth in floating point, almost
nothing in fixed point), and one extra state write (II~2, which halves
throughput but also halves the arithmetic). The closed loop buys a guarantee
that the encoder and decoder can never diverge, and II~2 buys a smaller design
the sensor does not notice; the float datapath buys only \SI{0.07}{\percent} of
compression ratio. What remains unexplained is the
replica's distance from Kiem's own LUT count, and Table~\ref{tab:kiemgap}
shows how much of it the testable explanations cover.
